\documentclass[11pt]{article}
\usepackage[margin=1in]{geometry}
\usepackage{booktabs}
\usepackage{tabularx}
\usepackage{array}
\usepackage{hyperref}
\usepackage{amsmath}

\newcolumntype{L}[1]{>{\raggedright\arraybackslash}p{#1}}

\title{Constraint Violation Rates in TabPFN:\\Preliminary Findings from an Initial Audit\\Across Relational and Tabular Benchmarks\\\large (with LightGBM as a Reference Baseline)}
\author{}
\date{}

\begin{document}
\maketitle

\begin{abstract}
This is an initial, preliminary audit of constraint-violation rate (CVR) in TabPFN, a genuine
tabular foundation model, evaluated on real relational and tabular benchmarks rather than
synthetic data. We ask whether TabPFN's predictions respect the integrity and domain
constraints that govern the databases it predicts over, using a constraint taxonomy spanning
eight generic database-integrity types (referential integrity, cardinality, aggregate
consistency, temporal order, functional dependency, denial constraint, range/domain,
monotonicity) plus domain-sourced rules requiring external knowledge. Every constraint is
verified against ground truth before being used to audit predictions, across three benchmarks:
rel-f1 (Formula 1 results), rel-salt (SAP enterprise sales data, 8 tasks), and HELOC (consumer
credit risk). TabPFN was run as a single, untuned configuration -- on rel-salt with a
necessarily small subsample (its CPU inference limits preclude full-scale training) and on
rel-f1 via RelArena's TabPFN-Rel relational pipeline -- so the TabPFN results reported here are
preliminary and should be read as a first pass, not a final characterization. LightGBM, a
conventional (non-foundation-model) gradient-boosted-tree baseline trained on the full
available data, is included throughout as a reference point, not as the paper's primary
subject, and its results are reported separately rather than merged into head-to-head
comparisons. In this initial pass, TabPFN's constraint-violation rates on rel-salt were low
across most tasks and it did not exhibit the LightGBM reference's one collapsed-prediction
failure mode; TabPFN was also more monotonic than the LightGBM reference on HELOC's
credit-risk features. We hold these as preliminary observations pending replication at matched
sample size and with a tuning sweep.
\end{abstract}

\section{Introduction}

A relational or tabular foundation model deployed to autocomplete a missing field, forecast an
outcome, or classify a row is implicitly making a claim about a structured world with its own
invariants: keys are unique, foreign keys resolve, aggregates sum correctly, some fields
determine others, and domain-specific rules hold. Accuracy against a held-out label says
nothing about whether a model's predictions, taken together, respect these invariants.

This report is an initial audit of one such model, TabPFN (a pretrained tabular foundation
model, not a model trained from scratch per task), for constraint-violation rate (CVR) across
real relational and tabular benchmarks. TabPFN is the paper's subject: it is the kind of model
the underlying research question is actually about, an architecture trained once and applied
to new tasks via in-context learning rather than gradient descent on each new dataset.
LightGBM, a conventional gradient-boosted-tree model trained from scratch on each task's own
full training data, is included as a familiar reference point for calibrating what a CVR
number means in absolute terms -- not because the paper is about comparing the two models, and
not on equal footing: LightGBM was trained on full datasets, while TabPFN's CPU inference
limits confined it to small subsamples on the largest benchmark. Results for the two models
are reported separately throughout, not merged into head-to-head tables, because they are not
run under comparable conditions.

We treat this as a preliminary, single-configuration pass: no hyperparameter tuning, no
repeated seeds, and (on rel-salt) a training sample far smaller than what LightGBM sees. Our
contributions are: (1) a constraint taxonomy spanning generic database-integrity properties and
domain-sourced rules, each with at least one real, empirically-verified instance; (2)
functional-dependency constraints derived from an enterprise database's own Operational
Business Knowledge Graph (OBKG) rather than hand-invented; and (3) an initial, preliminary
picture of TabPFN's constraint-violation behavior on real relational data, with LightGBM's
behavior reported alongside as reference context, including one case where a same-relation
consistency check alone would have understated a real failure in the reference baseline.

\section{Method}

\subsection{Constraint taxonomy}

We distinguish two tiers. \textbf{Tier 1 (database integrity)} covers eight generic constraint
types that hold for any well-formed relational database regardless of subject matter:
referential integrity, cardinality, aggregate consistency, temporal order, functional
dependency, denial constraint, range/domain, and monotonicity. \textbf{Tier 2
(domain-sourced)} covers rules whose justification requires external knowledge -- a sport's
rulebook, an ERP system's real business practice, or a fixed regulatory vocabulary -- rather
than following from the schema alone. The distinction is not always sharp (a functional
dependency can itself reflect real-world business knowledge), but it separates constraints
checkable from schema structure alone from those requiring a domain authority. Table
\ref{tab:examples} gives one real, verified example per type, drawn from this audit rather
than invented.

\begin{table}[h]
\centering
\small
\caption{One worked example per constraint type, with its measured ground-truth compliance.}
\label{tab:examples}
\begin{tabularx}{\linewidth}{@{} l L{4.3cm} X @{}}
\toprule
\textbf{Type} & \textbf{Example} & \textbf{Ground truth} \\
\midrule
Referential integrity & rel-f1: every \texttt{results.driverId} must exist in
  \texttt{drivers.driverId} & 0 / 26{,}080 violations \\[3pt]
Cardinality & rel-f1: at most one \texttt{results} row per (raceId, driverId) pair &
  85 / 25{,}989 violations (1950--78 shared drives) \\[3pt]
Aggregate consistency & rel-f1: a driver's final season standings points equal the sum of
  that driver's individual race points & 79 / 3{,}167 violations (no sprint-results table) \\[3pt]
Temporal order & rel-f1: a race's qualifying session date precedes the race date &
  0 / 9{,}815 violations \\[3pt]
Functional dependency & rel-salt: \texttt{SHIPPINGPOINT} determines \texttt{PLANT} (SAP
  config table TVSTZ: each shipping point belongs to one plant) & 1{,}356 / 2{,}319{,}540
  violations (99.94\% hold) \\[3pt]
Denial constraint & rel-salt: no two items in the same order may disagree on \texttt{PLANT}
  (ground truth) & 7{,}134 / 32{,}417{,}261 item pairs (99.98\% agree) \\[3pt]
Range / domain & rel-f1: a driver's points for a race must be non-negative &
  0 / 26{,}080 violations \\[3pt]
Monotonicity & HELOC: raising \texttt{ExternalRiskEstimate} (FICO's own risk score) must
  never lower predicted creditworthiness & real correlation $r=+0.46$ verified before use;
  violation rate is a prediction-level metric (Section \ref{sec:tabpfn-heloc}) \\
\bottomrule
\end{tabularx}
\end{table}

\subsection{Benchmarks}

\textbf{rel-f1}: RelBench's Formula 1 database, 1950--2023, three official tasks
(driver-position, driver-dnf, driver-top3). \textbf{rel-salt}: RelBench v2's SAP SALT-KG
enterprise sales database (2.3M line-item rows, 500K order headers), eight official
autocomplete tasks spanning item-level (item-plant, item-shippoint, item-incoterms) and
header-level (sales-office, sales-group, sales-payterms, sales-shipcond, sales-incoterms)
entities. \textbf{HELOC}: FICO's single-table consumer credit-risk dataset, audited for
monotonicity rather than relational constraints.

\subsection{Models}

\textbf{TabPFN v3 (primary model)}, run as a single untuned configuration in two different
integration paths depending on benchmark, both preliminary:

\begin{itemize}
\item On \textbf{rel-salt}, we used the plain \texttt{tabpfn} client package (v3 backend)
  directly on hand-built flat features, in-process. TabPFN's CPU inference has a documented
  10-class / $\sim$10{,}000-row soft advisory limit and a hard, unbypassable 160-class ceiling
  (confirmed empirically: sales-group's 502 classes triggered a \texttt{TabPFNValidationError}
  even with \texttt{ignore\_pretraining\_limits=True}, which we set to run past the soft
  10-class limit). Because rel-salt's full training tables run to 1.6M rows, we trained on an
  8{,}000-row random subsample and evaluated on a 1{,}000-row random subsample per task --
  small enough to fit TabPFN's practical CPU limits, but far smaller than what a from-scratch
  model like LightGBM can use, and not stratified by class (infeasible across up to 502
  classes in an 8{,}000-row budget). This is a feasibility trial, not a fully-powered
  evaluation.
\item On \textbf{rel-f1}, we used RelArena's \texttt{tabpfn-rel-local} pipeline (TabPFN-Rel),
  which automatically synthesizes relational features via multi-hop deep feature synthesis
  (DFS) across the database's foreign-key graph, rather than the hand-built flat features used
  on rel-salt, and performs its own internal tuning (10 trials) over a depth/hyperparameter
  grid. This is a materially different pipeline from the rel-salt integration, included
  because it is RelArena's standard way of running TabPFN on a relational benchmark; the two
  should not be read as the same experimental protocol.
\end{itemize}

\textbf{LightGBM (reference baseline)}, trained on the full available training data for every
task with a fixed, untuned configuration (200 estimators, 63 leaves) -- a conventional,
non-foundation-model baseline included so CVR numbers have a familiar point of reference, not
as the object of this study.

\subsection{Procedure}

Every constraint was first checked against real ground truth -- not assumed from a schema
diagram -- before being used to audit any model's predictions, following the same discipline
for all three benchmarks: compute the violation rate, and if it is unexpectedly non-zero,
either find and fix a real bug or document the non-zero rate as a genuine data property. This
discipline caught two implementation bugs during the audit (a temporal-censoring bug that made
both models degenerate into constant-class predictors, and a $\sim$1000x constraint-engine
performance bug from unprojected self-joins on rel-salt's 2.3M-row table), both fixed before
any result reported here was computed.

All decision-level violation rates reported in Section 3 are accompanied by a Wilson 95\%
confidence interval rather than a bare point estimate, since several checks -- most visibly
rel-salt's within-order consistency checks at $n=46$ item pairs -- rest on samples small
enough that a point estimate alone would overstate precision.

\paragraph{Scope of this initial pass.} This is a first, preliminary audit, not a finished
benchmark. Concretely: (1) TabPFN was run once per task, with no hyperparameter tuning on
rel-salt and RelArena's own fixed 10-trial tuning on rel-f1 -- neither is a tuning sweep
matched between the two models; (2) rel-salt's TabPFN subsample sizes (8{,}000 train /
1{,}000 test) were chosen for CPU feasibility, not statistical power -- the Wilson intervals
reported alongside each decision-level rate reflect that small sample, they do not remove its
effect; (3) each configuration was run once, not repeated across random seeds,
so we cannot yet distinguish a stable pattern from sampling noise; (4) LightGBM's numbers are
reported as a reference point throughout, never merged into a single head-to-head statistic
with TabPFN's, because the two were not run under comparable data-scale conditions. We draw
only preliminary conclusions from these results (Section \ref{sec:discussion}) and treat every
TabPFN number as provisional pending a matched-scale, repeated-seed follow-up.

\section{Results}

\subsection{Ground-truth verification}

Every constraint was checked against real data before being used to audit a model. Six of
eight rel-f1 constraints and five of eight rel-salt constraints hold exactly. This step is
model-independent: it validates the constraint specs themselves, not either model's
predictions.

\begin{table}[h]
\centering
\small
\caption{rel-f1 ground-truth constraint violations.}
\begin{tabular}{@{} l l r r r @{}}
\toprule
\textbf{Constraint} & \textbf{Type} & \textbf{Checked} & \textbf{Violations} & \textbf{Rate} \\
\midrule
results\_fk\_race & Referential integrity & 26{,}080 & 0 & 0.00\% \\
results\_fk\_driver & Referential integrity & 26{,}080 & 0 & 0.00\% \\
results\_fk\_constructor & Referential integrity & 26{,}080 & 0 & 0.00\% \\
results\_unique\_per\_race\_driver & Cardinality & 25{,}989 & 85 & 0.33\% \\
standings\_unique\_per\_race\_driver & Cardinality & 34{,}124 & 0 & 0.00\% \\
season\_points\_consistency & Aggregate consistency & 3{,}167 & 79 & 2.49\% \\
qualifying\_before\_race & Temporal order & 9{,}815 & 0 & 0.00\% \\
points\_non\_negative & Range/domain & 26{,}080 & 0 & 0.00\% \\
\bottomrule
\end{tabular}
\end{table}

The two non-zero rates are known, real data properties, not defects: the 85 duplicate (race,
driver) pairs are historical ``shared drives'' (1950--1978, two drivers legitimately sharing
one car); the 2.49\% aggregate mismatch arises because rel-f1 has no sprint-results table, so
2021+ sprint points are reflected in season standings but not recoverable by summing
individual race results.

\begin{table}[h]
\centering
\small
\caption{rel-salt ground-truth constraint violations.}
\begin{tabular}{@{} l l r r r @{}}
\toprule
\textbf{Constraint} & \textbf{Type} & \textbf{Checked} & \textbf{Violations} & \textbf{Rate} \\
\midrule
item\_fk\_salesdocument & Referential integrity & 2{,}319{,}540 & 0 & 0.00\% \\
item\_fk\_soldtoparty & Referential integrity & 2{,}319{,}540 & 0 & 0.00\% \\
customer\_fk\_address & Referential integrity & 139{,}607 & 0 & 0.00\% \\
order\_has\_at\_least\_one\_item & Cardinality & 500{,}908 & 0 & 0.00\% \\
salesorg\_determines\_division & Functional dependency & 500{,}908 & 0 & 0.00\% \\
salesorg\_determines\_billing\_company\_code & Functional dependency & 500{,}908 & 3 & 0.0006\% \\
shippingpoint\_determines\_plant & Functional dependency & 2{,}319{,}540 & 1{,}356 & 0.0585\% \\
salesgroup\_determines\_salesoffice & Functional dependency & 500{,}908 & 1{,}399 & 0.2793\% \\
\bottomrule
\end{tabular}
\end{table}

The four functional dependencies were derived from SAP SALT-KG's own Operational Business
Knowledge Graph, not hand-invented, and all four hold at or above 99.7\% -- consistent with
real enterprise master-data governance rather than a schema-declared key.

\subsection{TabPFN results (preliminary)}
\label{sec:tabpfn-results}

TabPFN's numbers below come from the single-configuration, small-sample runs described in
Sections 2.3--2.4; read them as an initial picture, not a final characterization.

\subsubsection{rel-salt: accuracy by task (1{,}000-row test sample)}

\begin{table}[h]
\centering
\small
\begin{tabular}{@{} l l r r @{}}
\toprule
\textbf{Task} & \textbf{Target} & \textbf{$n$} & \textbf{Accuracy} \\
\midrule
item-plant & PLANT & 1{,}000 & 98.90\% \\
item-shippoint & SHIPPINGPOINT & 1{,}000 & 96.90\% \\
item-incoterms & ITEMINCOTERMSCLASSIFICATION & 1{,}000 & 100.00\% \\
sales-office & SALESOFFICE & 1{,}000 & 100.00\% \\
sales-group & SALESGROUP & -- & excluded (502 classes $>$ 160-class ceiling) \\
sales-payterms & CUSTOMERPAYMENTTERMS & 1{,}000 & 98.90\% \\
sales-shipcond & SHIPPINGCONDITION & 1{,}000 & 98.80\% \\
sales-incoterms & HEADERINCOTERMSCLASSIFICATION & 1{,}000 & 100.00\% \\
\bottomrule
\end{tabular}
\end{table}

\subsubsection{rel-salt: prediction-level constraint violations}

\begin{table}[h]
\centering
\small
\begin{tabularx}{\linewidth}{@{} L{5.7cm} l r l @{}}
\toprule
\textbf{Constraint} & \textbf{$n$} & \textbf{TabPFN CVR} & \textbf{95\% CI (Wilson)} \\
\midrule
plant\_consistent\_within\_order & 46 item pairs & 0.00\% (0) & [0.00, 7.71] \\
shippingpoint\_consistent\_within\_order & 46 item pairs & 0.00\% (0) & [0.00, 7.71] \\
itemincotermsclassification\_consistent\_within\_order & 46 item pairs & 0.00\% (0) & [0.00, 7.71] \\
shippingpoint\_pred\_determines\_plant\_pred & 1{,}000 & 0.50\% (5) & [0.21, 1.17] \\
salesgroup\_pred\_determines\_salesoffice\_pred & -- & excluded (sales-group not runnable) & -- \\
\bottomrule
\end{tabularx}
\end{table}

\subsubsection{rel-f1 (TabPFN-Rel via RelArena): prediction-level constraint violations}

\begin{table}[h]
\centering
\small
\begin{tabular}{@{} l l r l @{}}
\toprule
\textbf{Constraint} & \textbf{$n$} & \textbf{TabPFN-Rel CVR} & \textbf{95\% CI (Wilson)} \\
\midrule
top3\_predictions\_per\_window & 30 windows & 36.67\% (11) & [21.87, 54.49] \\
position\_pred\_in\_domain & 760 & 0.00\% (0) & [0.00, 0.50] \\
\bottomrule
\end{tabular}
\end{table}

\subsubsection{HELOC: monotonicity}
\label{sec:tabpfn-heloc}

\begin{table}[h]
\centering
\small
\begin{tabular}{@{} l l r l @{}}
\toprule
\textbf{Feature} & \textbf{$n$} & \textbf{TabPFN CVR} & \textbf{95\% CI (Wilson)} \\
\midrule
ExternalRiskEstimate & 1{,}984 & 10.53\% (209) & [9.26, 11.96] \\
NetFractionRevolvingBurden & 1{,}945 & 19.69\% (383) & [17.98, 21.52] \\
\bottomrule
\end{tabular}
\end{table}

In this initial pass, TabPFN's rel-salt accuracy was high ($\geq$96.9\%) on every task it
could run, its prediction-level constraint violations were at or below 0.5\% everywhere
measured, and its rel-f1 top3\_predictions\_per\_window rate (36.67\%, 11 of 30 windows) is
the one figure worth flagging for a closer look given how few test windows back it.

\subsection{LightGBM reference results}
\label{sec:lgbm-results}

LightGBM was trained on the full available data for every task (Section 2.3) and is reported
here only as a reference point for what these CVR numbers look like for a conventional,
fully-trained baseline -- not as a comparison target for TabPFN.

\subsubsection{rel-salt: accuracy by task (full test set)}
\label{sec:lgbm-accuracy}

\begin{table}[h]
\centering
\small
\begin{tabular}{@{} l l r r @{}}
\toprule
\textbf{Task} & \textbf{Target} & \textbf{$n$} & \textbf{Accuracy} \\
\midrule
item-plant & PLANT & 402{,}855 & 99.08\% \\
item-shippoint & SHIPPINGPOINT & 402{,}855 & 2.97\% \\
item-incoterms & ITEMINCOTERMSCLASSIFICATION & 402{,}855 & 99.99\% \\
sales-office & SALESOFFICE & 88{,}942 & 100.00\% \\
sales-group & SALESGROUP & 88{,}942 & 86.04\% \\
sales-payterms & CUSTOMERPAYMENTTERMS & 88{,}942 & 99.85\% \\
sales-shipcond & SHIPPINGCONDITION & 88{,}942 & 18.12\% \\
sales-incoterms & HEADERINCOTERMSCLASSIFICATION & 88{,}942 & 99.98\% \\
\bottomrule
\end{tabular}
\end{table}

LightGBM's item-shippoint accuracy (2.97\%) falls below the 48.2\% majority-class baseline
for that task; sales-shipcond shows the same pattern (18.12\%). We do not have a confirmed
explanation (Section \ref{sec:limitations}).

\subsubsection{rel-salt: prediction-level constraint violations}

\begin{table}[h]
\centering
\small
\begin{tabularx}{\linewidth}{@{} L{5.7cm} l r l @{}}
\toprule
\textbf{Constraint} & \textbf{$n$} & \textbf{LightGBM CVR} & \textbf{95\% CI (Wilson)} \\
\midrule
plant\_consistent\_within\_order & 6{,}184{,}238 item pairs & 0.0078\% (483) & [0.0071, 0.0085] \\
shippingpoint\_consistent\_within\_order & 6{,}184{,}238 item pairs & 0.00\% (0) & [0.00, 0.0001] \\
itemincotermsclassification\_consistent\_within\_order & 6{,}184{,}238 item pairs & 0.0049\% (303) & [0.0044, 0.0055] \\
shippingpoint\_pred\_determines\_plant\_pred & 402{,}855 & 27.29\% (109{,}941) & [27.15, 27.43] \\
salesgroup\_pred\_determines\_salesoffice\_pred & 88{,}942 & 0.11\% (102) & [0.09, 0.14] \\
\bottomrule
\end{tabularx}
\end{table}

The shippingpoint\_consistent\_within\_order row (0.00\%) and the
shippingpoint\_pred\_determines\_plant\_pred row (27.29\%) both concern the same
item-shippoint predictions; Section \ref{sec:masking} discusses why they disagree so sharply.

\subsubsection{rel-f1 (LightGBM): prediction-level constraint violations}

\begin{table}[h]
\centering
\small
\begin{tabular}{@{} l l r l @{}}
\toprule
\textbf{Constraint} & \textbf{$n$} & \textbf{LightGBM CVR} & \textbf{95\% CI (Wilson)} \\
\midrule
top3\_predictions\_per\_window & 30 windows & 13.33\% (4) & [5.31, 29.68] \\
position\_pred\_in\_domain & 760 & 0.00\% (0) & [0.00, 0.50] \\
\bottomrule
\end{tabular}
\end{table}

\subsubsection{HELOC: monotonicity}

\begin{table}[h]
\centering
\small
\begin{tabular}{@{} l l r l @{}}
\toprule
\textbf{Feature} & \textbf{$n$} & \textbf{LightGBM CVR} & \textbf{95\% CI (Wilson)} \\
\midrule
ExternalRiskEstimate & 1{,}984 & 17.84\% (354) & [16.22, 19.59] \\
NetFractionRevolvingBurden & 1{,}945 & 32.19\% (626) & [30.15, 34.29] \\
\bottomrule
\end{tabular}
\end{table}

\section{Discussion}
\label{sec:discussion}

\subsection{TabPFN's preliminary CVR pattern}

In this initial, single-configuration pass, TabPFN's rel-salt predictions showed low
constraint-violation rates everywhere we could measure them (at or below 0.5\%) alongside
high accuracy ($\geq$96.9\% on every runnable task), and it did not reproduce the LightGBM
reference's item-shippoint collapse (Section \ref{sec:masking}). TabPFN was also more
monotonic than the LightGBM reference on both HELOC features (10.53\%/19.69\% vs.\
17.84\%/32.19\%). We read these as encouraging preliminary signals, not a settled finding: the
rel-salt sample is small (1{,}000 test rows per task), untuned, and run once, so we cannot yet
separate a real property of TabPFN's predictions from sampling variation. The one number that
stands out enough to flag for follow-up is rel-f1's top3\_predictions\_per\_window via
TabPFN-Rel (36.67\%, 11 of 30 windows) -- notably higher than the LightGBM reference's figure
on the same check (13.33\%, 4 of 30), though with only 30 test windows a handful of counts
move the percentage a great deal: a two-sided Fisher exact test on the 11-vs-4-of-30 counts
gives $p=0.07$, so at this sample size the gap does not reach significance. TabPFN-Rel's
RelArena pipeline is also a different feature-engineering path from anything used on rel-salt
(Section 2.3). We therefore read this result as evidence that per-cell predictors generally
lack a mechanism for enforcing this kind of cross-row cardinality constraint, not as a ranking
between the two models audited here, and not as evidence about TabPFN specifically.

\subsection{Why the LightGBM reference needed more than one constraint}
\label{sec:masking}

One result from the LightGBM reference baseline illustrates why this audit uses more than one
constraint type per relation, and we include it because it shapes how the TabPFN numbers above
should be read too. LightGBM's within-order consistency check reports 0.00\% violations for
its item-shippoint predictions (Section \ref{sec:lgbm-results}) -- an apparently perfect
result. But item-shippoint's LightGBM accuracy is 2.97\%, below the task's 48.2\%
majority-class baseline (Section \ref{sec:lgbm-accuracy}): the predictions are not just wrong,
they concentrate on a handful of classes, so items in the same order trivially agree with each
other on a wrong answer. The cross-task functional-dependency check
(shippingpoint\_pred\_determines\_plant\_pred, derived from the same OBKG relation used to
verify the ground-truth \texttt{SHIPPINGPOINT}$\to$\texttt{PLANT} dependency) exposes the same
failure directly, at 27.29\%, because a model whose shippingpoint predictions are near-random
cannot also produce a plant prediction consistent with them.

We take two things from this. First, methodologically: a single same-relation consistency
check is not sufficient evidence of model competence, for either model audited here -- it is
why TabPFN's results above are reported across several independent constraint types rather
than one. Second, and more tentatively: TabPFN did not show this failure mode on the same task
(96.90\% accuracy, Section \ref{sec:tabpfn-results}), which is worth a matched-scale follow-up
rather than a conclusion, since the two models were not run on equal-sized samples.

\subsection{Constraint engineering at real scale}

A naive self-join implementation of the within-order and functional-dependency checks took
15--65 minutes per constraint on rel-salt's 2.3M-row item table, driven by nullable-dtype
columns carried unnecessarily through the join. Projecting each join down to only the columns
the predicate needs reduced this to under a second per constraint with identical violation
counts. Constraint auditing at production scale is not just a modeling exercise; the
constraint engine itself must be engineered for the data volumes real relational benchmarks
involve.

\section{Limitations}
\label{sec:limitations}

\begin{itemize}
\item \textbf{Unequal statistical power.} TabPFN's rel-salt results use an 8{,}000-row
  training subsample and a 1{,}000-row test subsample against LightGBM's full $\sim$1.6M
  training / $\sim$400K test rows, so the two models' violation rates are not equal-power
  comparisons. On item-level tasks, TabPFN's feature set also drops five near-unique
  identifier columns that LightGBM retains (Section 2.3).
\item \textbf{Single run, no matched tuning.} TabPFN was run once per task, untuned on
  rel-salt, and under RelArena's own fixed 10-trial tuning on rel-f1 -- a different,
  non-matched regime. We cannot yet distinguish a stable pattern from a lucky or unlucky draw.
\item \textbf{A hard architectural ceiling.} TabPFN has an unbypassable 160-class limit that
  excludes it entirely from sales-group (502 classes); every comparison involving that task is
  one-sided.
\item \textbf{An unexplained weak point in the LightGBM reference.} LightGBM's item-shippoint
  and sales-shipcond weakness (Section \ref{sec:lgbm-accuracy}) has a plausible explanation --
  distribution shift toward customer identities unseen at training time -- that we did not
  test; raising \texttt{n\_estimators} from 200 to 800 did not change the result, ruling out
  simple undertraining. This is a property of the reference baseline, not of TabPFN.
\item \textbf{Two different TabPFN pipelines.} rel-f1's TabPFN-Rel (RelArena, DFS relational
  features, internal tuning) and rel-salt's plain TabPFN (hand-built flat features, no
  tuning) are not the same experimental setup; results from the two should not be pooled or
  treated as evidence about a single ``TabPFN'' condition.
\item \textbf{Small denominators.} rel-f1's top3\_predictions\_per\_window check has only 30
  test windows; an 11-vs-4 count difference reads as a large percentage gap (36.7\% vs.\
  13.3\%).
\end{itemize}

\section{Conclusion (preliminary)}

This is an initial, single-configuration audit, not a finished evaluation. In this first pass,
TabPFN showed low constraint-violation rates and high accuracy across most of rel-salt's
tasks, and did not reproduce the one collapsed-prediction failure mode seen in the LightGBM
reference baseline; it was also more monotonic than that reference on HELOC. We hold all of
these as preliminary observations: TabPFN's rel-salt runs used a small, untuned, single-seed
sample, and its rel-f1 numbers come from a different pipeline (TabPFN-Rel) than its rel-salt
numbers, so the two should not be pooled into one claim about ``TabPFN.'' The audit's more
durable finding is methodological: a same-relation consistency check gave a misleadingly
perfect result for the LightGBM reference's one collapsed-prediction task, and only an
independently-sourced, cross-task constraint exposed it -- a caution that applies to auditing
any model, including TabPFN, going forward. A matched-sample-size, repeated-seed, tuned
follow-up is the natural next step before drawing firmer conclusions about TabPFN's constraint
compliance.

\end{document}
